\section{Miaodong Technology: metodología de emprendimiento}

\subsection{La lógica de la sede en dos ciudades}

Miaodong Technology estableció una oficina en Shenzhen y otra en Silicon Valley. Yang Shuo (杨硕) considera que esta sede en dos ciudades es necesaria:

\begin{itemize}
    \item \textbf{Ventajas de Silicon Valley}: una cultura que no teme al fracaso, que fomenta la innovación, con movilidad activa de personal, profesionales de pensamiento ágil y una exploración activa de nuevas tecnologías
    \item \textbf{Ventajas de China}: un sistema de cadena de suministro maduro, que es una ventaja mucho mayor que la de Silicon Valley; el país también fomenta los avances tecnológicos y la innovación
    \item \textbf{Necesidad de integración}: integrar en un solo equipo a personas de distintos orígenes y con ideas distintas, para que cada quien aproveche sus propias fortalezas
\end{itemize}

\subsection{La ruta 2C y el ciclo comercial cerrado}

Miaodong Technology eligió una ruta \textbf{orientada al 2C} (dirigida al consumidor). El razonamiento de Yang Shuo es el siguiente:

\begin{enumerate}
    \item Primero encontrar un escenario de aplicación concreto
    \item Elaborar la prueba de concepto y venderla
    \item Recorrer por completo el proceso de producción en volumen, posventa y reparación
    \item Sobre esa base, ampliar de manera gradual la capacidad de inteligencia
\end{enumerate}

Él enfatiza en particular que, en materia de inteligencia, adoptarán un enfoque más abierto y una estrategia orientada al ecosistema. En cuanto a tecnologías de Foundation Model como VLA (Vision-Language-Action), aunque son muy prometedoras, la cantidad de recursos que exigen es comparable a la de desarrollar hardware. En la etapa actual, muy pocas empresas que obtienen financiamiento puramente de inversión financiera logran hacer bien, al mismo tiempo, tanto el hardware como el software: «es posible que al final ambas cosas queden hechas al 60\,\%». Por eso, la mejor forma de organización es que distintas empresas se dividan el trabajo y colaboren, concentrándose cada una en un aspecto.

\begin{knowledgebox}{La división del trabajo entre las empresas de Foundation Model y las empresas de hardware}
Yang Shuo observa que en la industria de la robótica existen dos tipos de empresas emergentes: unas que primero encuentran un escenario y elaboran el producto (como Miaodong Technology), y otras que primero desarrollan tecnología de Foundation Model (como VLA e Imitation Learning). Él considera que ambas son necesarias: son demasiados los problemas que hay que superar en todo este proceso, y se necesitan equipos distintos que exploren por separado, encontrando cada uno, mediante la colaboración y el uso compartido de recursos, una forma comercialmente sostenible de desarrollarse. Desde luego, la situación ideal sería que una sola empresa lo hiciera todo junto (como Tesla), pero eso exige una inversión de recursos demasiado grande para una empresa emergente.
\end{knowledgebox}

\subsection{El lado positivo de la competencia}

En 2025, la competencia entre las empresas emergentes de robótica es «muy intensa»: hay una enorme competencia en todos los frentes, desde la inversión y la contratación de personal hasta los proveedores y la tecnología de los modelos. Pero Yang Shuo considera que esto es positivo:

\begin{itemize}
    \item Los inversionistas se dan cuenta de que la industria terminará por producir productos exitosos
    \item La captación y formación de talento representa una contribución de largo plazo para la industria en su conjunto
    \item Si se logra encontrar un escenario comercial que permita la producción en volumen a gran escala de componentes, esto también impulsa la cadena de suministro
    \item La competencia sana ayuda a mantener una mentalidad de avance rápido
\end{itemize}

\subsection{La convergencia tecnológica y el verdadero foso competitivo}

Yang Shuo observa que las rutas tecnológicas de la industria son cada vez más convergentes: el hardware se parece cada vez más entre sí, y lo mismo ocurre con el Foundation Model. Él considera que la convergencia tecnológica es normal; una tecnología madura converge inevitablemente hacia la solución de mayor eficiencia y mejor relación costo-beneficio, tal como el automóvil terminó por converger hacia ciertas líneas de diseño comunes.

Entonces, ¿dónde está el verdadero foso competitivo de una empresa emergente?

\begin{importantbox}{El foso competitivo de las empresas emergentes: los detalles y la gestión organizacional}
Yang Shuo considera que, en una época en la que las barreras tecnológicas son cada vez más superficiales (el aprendizaje profundo hizo que muchas tecnologías que antes podían «mantenerse ocultas» fueran desplazadas por métodos nuevos, y además la tecnología se comparte en el ámbito académico), el foso competitivo de las empresas emergentes ha regresado a la \textbf{gestión organizacional} y a los \textbf{detalles}: los detalles del entrenamiento de modelos, los detalles del despliegue, los detalles del hardware y del software; la acumulación de estos detalles produce enormes diferencias. Tomando a DJI como ejemplo: las soluciones de control de vuelo, cámara y transmisión de imagen ya están maduras desde hace tiempo en la industria, pero DJI logra hacer un producto mejor que los demás precisamente gracias al cuidado extremo que pone en cada eslabón técnico y en cada detalle de la interacción del producto.
\end{importantbox}

\subsection{Si yo fuera inversionista}

Al preguntársele «si fueras inversionista, ¿en qué tipo de empresa de robótica invertirías?», Yang Shuo respondió que prestaría atención a tecnologías de base más a largo plazo, por ejemplo, nuevas soluciones de motores y nuevos materiales que permitan a los robots trabajar de manera continua en entornos de alta temperatura. Aunque esta tecnología no se use en algún robot humanoide concreto dentro de cinco años, a largo plazo sin duda tendrá valor. Por supuesto, también exigiría que la empresa tuviera capacidad de comercialización a corto plazo (por ejemplo, vender motores de articulación a los fabricantes actuales de robots humanoides), para lograr un desarrollo sostenible.

En cuanto a las perspectivas de mercado, Yang Shuo está convencido de que el robot humanoide es un mercado del orden de cientos de miles de millones a billones. Una vez que la tecnología madure, habrá en el mundo entre 10 y 15 grandes empresas de robótica, tal como ocurre hoy en el mercado automotriz: distintas empresas con distintos posicionamientos (como el sedán, el auto deportivo, el vehículo comercial), y los consumidores eligen entre varias marcas.

\subsection{Resumen de la sección}

La metodología de emprendimiento de Miaodong Technology puede resumirse así: una sede en dos ciudades que aprovecha las ventajas de ambos lugares, una ruta 2C que prioriza recorrer por completo toda la cadena del producto, una postura abierta pero no de acaparamiento directo frente a tecnologías de vanguardia como VLA, y la construcción de un foso competitivo mediante la búsqueda de detalle y la gestión organizacional. En un contexto de competencia intensa en la industria, mantiene una visión de largo plazo y confía en que el robot humanoide terminará por convertirse en un mercado de billones.

